\section{A sharp Gram-map criterion for HHC}
\label{sec:gram-hyperplanes}

We now turn from certificates to exact hull descriptions. This section
gives a structural HHC criterion for any number of forms with repeated
quadratic blocks and no linear terms. The matrix variational identity in
its proof is classical; the conclusion concerns all real hyperplane
restrictions of the full Gram map.

For positive integers $k,r$, consider
\[
 \Phi_{k,r}:\R^{k\times r}\times\R\longrightarrow\Sym^k\times\R,
 \qquad \Phi_{k,r}(X,t)=(XX\transpose,t^2).
\]
We use the Frobenius inner product
$\langle B,X\rangle=\tr(B\transpose X)$ for rectangular matrices.
For $H,G\succeq0$, define
\begin{equation}\label{eq:squared-fidelity}
 F_H(G)=\left[\tr\left((G^{1/2}HG^{1/2})^{1/2}\right)\right]^2.
\end{equation}
This is the squared fidelity, also called transition probability, of
the two positive semidefinite matrices, without a trace normalization.

\begin{theorem}[Sharp Gram-map threshold]\label{thm:gram-hhc}
The map $\Phi_{k,r}$ has HHC if and only if $r\geq k$.
If $r\geq k$ and $(k,r)\ne(1,1)$, then for every hyperplane
\[
 L=\{(X,t):\langle B,X\rangle+\beta t=0\},
 \qquad (B,\beta)\ne(0,0),
\]
its image is exactly
\begin{equation}\label{eq:gram-hyperplane-image}
 \Phi_{k,r}(L)=
 \{(G,T):G\succeq0,\ T\geq0,\ \beta^2T\leq F_{BB\transpose}(G)\}.
\end{equation}
The formula includes singular $G$ or $BB\transpose$ and zero values of
either hyperplane coefficient.
\end{theorem}

We prove the attainment and concavity statements separately. The
following real trace-range argument handles the disconnectedness of the
orthogonal group by an elementary rotation construction. Singular-value
optimization over rotations is classical; see, for example,
\citet[Lemma~1]{RamachandranShuWang2025}, which credits the
determinant-corrected $SO(k)$ extremum to earlier work.

\begin{lemma}[Linear functionals on a Gram fiber]\label{lem:gram-range}
Let $r\geq k$, $G\succeq0$, and $B\in\R^{k\times r}$. The maximum
of $\langle B,X\rangle$ over $XX\transpose=G$ equals
\[
 M=\tr\left((G^{1/2}BB\transpose G^{1/2})^{1/2}\right).
\]
If $(k,r)\ne(1,1)$, the attained values form the entire interval
$[-M,M]$.
\end{lemma}
\begin{proof}
Every factor of $G$ can be written as $X=G^{1/2}Q$ with
$QQ\transpose=I_k$. Indeed, in an eigenbasis of $G$, normalize the
rows belonging to positive eigenvalues and extend those orthonormal
rows to $k$ orthonormal rows in $\R^r$. Rows belonging to zero
eigenvalues of $G$ were already zero, so this extension preserves $X$.

Take a rectangular singular value decomposition
$G^{1/2}B=U[D\ 0]V\transpose$, where
$D=\operatorname{diag}(s_1,\ldots,s_k)$ and
$s_1\geq\cdots\geq s_k\geq0$. Then
\[
 \langle B,X\rangle
 =\sum_{j=1}^k s_j(U\transpose QV)_{jj}
 \leq\sum_{j=1}^k s_j=M,
\]
because every displayed diagonal entry has absolute value at most one.
The choice $Q=U[I_k\ 0]V\transpose$ attains $M$; its negative
attains $-M$.

For $k\geq2$, restrict to $Q=U[R\ 0]V\transpose$, with $R$
orthogonal. If $k$ is even, rotate each disjoint pair of coordinate
axes through an angle $\theta\in[0,\pi]$. This gives a continuous
path from $I_k$ to $-I_k$, along which $\tr(DR)$ runs continuously
from $M$ to $-M$. If $k$ is odd, then $k\geq3$. Leave the last axis
fixed and rotate the other pairs. The endpoint value is
$2s_k-M\leq0$, since $M\geq ks_k$. The intermediate value theorem
therefore gives every value in $[0,M]$. Negating the factors gives
$[-M,0]$. If $k=1$ and $r\geq2$, the nonzero Gram fiber is a sphere
in $\R^r$; a planar rotation from a maximizing vector to its negative
gives the same conclusion. The zero fiber causes no exception.
\end{proof}

\begin{lemma}[Classical squared-fidelity formula]\label{lem:fidelity-concavity}
For $G,H\succeq0$,
\begin{equation}\label{eq:fidelity-infimum}
 F_H(G)=\inf_{Z\succ0}\tr(GZ)\tr(HZ^{-1}).
\end{equation}
Consequently $G\mapsto F_H(G)$ is concave on $\Sym^k_+$ for fixed $H$.
\end{lemma}

Formula~\eqref{eq:fidelity-infimum} and separate concavity are prior
results in fidelity theory: see \citet[pp.~408--410, equation~(14) with
the partial-fidelity index equal to zero]{Uhlmann2000}. We include a
real-matrix proof to account explicitly for singular matrices and the
absence of trace normalization.

\begin{proof}
Use Lemma~\ref{lem:gram-range} with $r=k$ and $B=H^{1/2}$ to choose
$XX\transpose=G$ with $\langle B,X\rangle=\sqrt{F_H(G)}$.
For every $Z\succ0$, the Frobenius Cauchy--Schwarz inequality gives
\[
 F_H(G)=\langle Z^{-1/2}B,Z^{1/2}X\rangle^2
 \leq\tr(GZ)\tr(HZ^{-1}).
\]
If $G,H\succ0$, put $K=(G^{1/2}HG^{1/2})^{1/2}$ and
$Z=G^{-1/2}KG^{-1/2}$. Then $ZGZ=H$, and
$\tr(GZ)=\tr(HZ^{-1})=\tr K$, so equality holds.

For general positive semidefinite $G,H$, apply this construction to
$G_\epsilon=G+\epsilon I$ and $H_\epsilon=H+\epsilon I$, obtaining
$Z_\epsilon\succ0$. Positivity gives
\[
 \inf_{Z\succ0}\tr(GZ)\tr(HZ^{-1})
 \leq\tr(GZ_\epsilon)\tr(HZ_\epsilon^{-1})
 \leq F_{H_\epsilon}(G_\epsilon).
\]
Continuity of positive semidefinite matrix square roots lets
$\epsilon\downarrow0$ and proves the reverse inequality.
For fixed $H$, each expression in the infimum is linear in $G$;
an infimum of linear functions is concave.
\end{proof}

\begin{proof}[Proof of Theorem~\ref{thm:gram-hhc}]
Suppose $r\geq k$ and $(k,r)\ne(1,1)$. Given $G\succeq0$ and
$T\geq0$, take $t=\sqrt T$. Lemma~\ref{lem:gram-range} says that
one can satisfy $\langle B,X\rangle=-\beta t$ with $XX\transpose=G$
if and only if $|\beta|\sqrt T\leq\sqrt{F_{BB\transpose}(G)}$.
This proves~\eqref{eq:gram-hyperplane-image}. Its right-hand side
is convex by Lemma~\ref{lem:fidelity-concavity}, proving HHC.
When $k=r=1$, every hyperplane in the two-dimensional domain is a
line. Its quadratic image is a ray, hence convex. In that case the
image formula need not hold: $x+t=0$ has image
$\{(G,T):G=T\geq0\}$, rather than $0\leq T\leq G$.

Conversely, if $r<k$, the image of the hyperplane $t=0$ is
$\{(G,0):G\succeq0,\ \operatorname{rank}G\leq r\}$.
The matrices $G_1=\sum_{j=1}^r e_je_j\transpose$ and
$G_2=e_{r+1}e_{r+1}\transpose$ belong to this image, but their
midpoint has rank $r+1$. The image is not convex.
\end{proof}

\begin{corollary}[Repeated quadratic blocks]\label{cor:gram-family}
If $r\geq k$, then any finite family
\[
 q_i(X,t)=\tr(A_iXX\transpose)+c_it^2,
 \qquad A_i\in\Sym^k,\quad c_i\in\R,
\]
has HHC, with no restriction on the number of forms.
\end{corollary}
\begin{proof}
On every hyperplane the image of $(q_1,\ldots,q_m)$ is a linear
image of the convex set $\Phi_{k,r}(L)$.
\end{proof}

The necessity of $r\geq k$ concerns the \emph{full} Gram map; a
smaller collection of its output coordinates can have HHC below this
threshold. Nor does the corollary allow arbitrary mixed terms
$2t\langle D_i,X\rangle$. For $k=2$ the image formula simplifies to
\[
 F_H(G)=\tr(HG)+2\sqrt{\det H\det G},
\]
by expanding the square of the sum of two singular values. This
two-row case yields the four-variable construction below.

The classical ingredients are the singular-value extremum and
squared-fidelity concavity; the conclusion here is their sharp HHC
formulation for a Gram map with one additional scalar square. Results
on quadratic matrix programming and rotation images provide related
convexity mechanisms
\citep{Beck2007,RamachandranShuWang2025}. In particular,
\citet[Theorems~3.1 and~3.4]{Beck2009} proves global image convexity
for real quadratic matrix functions when their number is at most the
replication count, and for one additional function under a positive
definite combination of the quadratic blocks and a row dimension of
at least two. These are global image results; an arbitrary hyperplane
restriction generally destroys that repeated-block structure.
Theorem~\ref{thm:gram-hhc} instead preserves every Gram entry and
the scalar square while imposing one arbitrary real linear equation.
The next section explains the SDP comparison relevant to its application.
